\documentclass[11pt, a4paper]{article}
\usepackage[a4paper, top=2.5cm, bottom=2.5cm, left=2cm, right=2cm]{geometry}
\usepackage{fontspec}
\usepackage[english, bidi=basic, provide=*]{babel}
\babelprovide[import, onchar=ids fonts]{english}
\babelfont{rm}{TeX Gyre Termes}
\babelfont{sf}{TeX Gyre Heros}
\babelfont{tt}{TeX Gyre Cursor}
\usepackage{enumitem}
\usepackage{amsmath}

\title{MCQ Practice Set: A Tour of the Cell}
\author{Subject: Biology}
\date{}

\begin{document}
\maketitle

\section*{Practice Questions}

\begin{enumerate}[leftmargin=*, label=\arabic*.]

\item \textbf{Which parameter of microscopy refers to the ratio of an object's image size to its real size?}
\begin{enumerate}[label=\Alph*.]
\item Resolution
\item Contrast
\item Magnification
\item Indexing
\item Refraction
\end{enumerate}
\textbf{Answer:} C \\
\textbf{Explanation:} Magnification is the ratio of an object's image size to its actual size.

\item \textbf{A biologist wants to study the 3D surface topography of a cell. Which tool is most appropriate?}
\begin{enumerate}[label=\Alph*.]
\item Light Microscope (LM)
\item Transmission Electron Microscope (TEM)
\item Scanning Electron Microscope (SEM)
\item Phase-Contrast Microscope
\item Super-resolution Microscope
\end{enumerate}
\textbf{Answer:} C \\
\textbf{Explanation:} SEM scans the surface of a specimen, usually coated in gold, to create a 3D-appearing image of the surface.

\item \textbf{What is the theoretical resolution limit of a standard light microscope?}
\begin{enumerate}[label=\Alph*.]
\item 2 nm
\item 200 nm
\item 2 $\mu$m
\item 20 nm
\item 0.2 nm
\end{enumerate}
\textbf{Answer:} B \\
\textbf{Explanation:} Standard light microscopy cannot resolve detail finer than about 0.2 $\mu$m or 200 nm.

\item \textbf{Which type of microscopy uses lasers to provide a single plane of fluorescence, eliminating out-of-focus light?}
\begin{enumerate}[label=\Alph*.]
\item Brightfield
\item Phase-contrast
\item Differential interference contrast
\item Confocal
\item Deconvolution
\end{enumerate}
\textbf{Answer:} D \\
\textbf{Explanation:} Confocal microscopy uses a laser to scan a single plane, providing sharper 3D reconstructions.

\item \textbf{In cell fractionation, which organelle would typically settle into the pellet at the lowest centrifugation speed?}
\begin{enumerate}[label=\Alph*.]
\item Ribosomes
\item Mitochondria
\item Chloroplasts
\item Nuclei
\item Microsomes
\end{enumerate}
\textbf{Answer:} D \\
\textbf{Explanation:} Nuclei are the largest and densest components, so they settle first at 1,000 g.

\item \textbf{What is the primary advantage of light microscopy over electron microscopy?}
\begin{enumerate}[label=\Alph*.]
\item Higher resolution
\item Ability to observe living cells
\item Greater magnification
\item Visualization of ribosomes
\item Use of heavy metal stains
\end{enumerate}
\textbf{Answer:} B \\
\textbf{Explanation:} Electron microscopy preparation usually kills the cells, whereas light microscopy allows the study of living processes.

\item \textbf{Which of the following is found in both prokaryotic and eukaryotic cells?}
\begin{enumerate}[label=\Alph*.]
\item Mitochondria
\item Golgi apparatus
\item Ribosomes
\item Endoplasmic reticulum
\item Nuclear envelope
\end{enumerate}
\textbf{Answer:} C \\
\textbf{Explanation:} Both cell types contain ribosomes for protein synthesis, though their sizes differ.

\item \textbf{The region where DNA is concentrated in a prokaryotic cell is called the:}
\begin{enumerate}[label=\Alph*.]
\item Nucleus
\item Nucleolus
\item Nucleoid
\item Nucleosome
\item Cytoplasm
\end{enumerate}
\textbf{Answer:} C \\
\textbf{Explanation:} Prokaryotic DNA is concentrated in the nucleoid, a region not enclosed by a membrane.

\item \textbf{As a cell increases in size, the \underline{\hspace{1cm}} grows proportionately more than its \underline{\hspace{1cm}}.}
\begin{enumerate}[label=\Alph*.]
\item Surface area; volume
\item Volume; surface area
\item Length; width
\item Radius; diameter
\item Mass; density
\end{enumerate}
\textbf{Answer:} B \\
\textbf{Explanation:} Volume increases as the cube of the radius ($r^3$), while surface area increases as the square ($r^2$).

\item \textbf{Which structure is responsible for maintaining the shape of the nucleus?}
\begin{enumerate}[label=\Alph*.]
\item Nuclear pore
\item Nuclear lamina
\item Nucleolus
\item Chromatin
\item Endoplasmic reticulum
\end{enumerate}
\textbf{Answer:} B \\
\textbf{Explanation:} The nuclear lamina is a netlike array of protein filaments that lines the inner surface of the envelope.

\item \textbf{What is the site of ribosomal RNA (rRNA) synthesis in eukaryotes?}
\begin{enumerate}[label=\Alph*.]
\item Cytosol
\item Nucleolus
\item Smooth ER
\item Rough ER
\item Golgi apparatus
\end{enumerate}
\textbf{Answer:} B \\
\textbf{Explanation:} The nucleolus is the region in the nucleus where rRNA is synthesized and ribosomal subunits are assembled.

\item \textbf{Which organelle is continuous with the nuclear envelope?}
\begin{enumerate}[label=\Alph*.]
\item Golgi apparatus
\item Lysosome
\item Mitochondrion
\item Endoplasmic reticulum
\item Peroxisome
\end{enumerate}
\textbf{Answer:} D \\
\textbf{Explanation:} The ER membrane is physically continuous with the outer membrane of the nuclear envelope.

\item \textbf{A cell with a high rate of detoxification, such as a liver cell, would be expected to have an abundance of:}
\begin{enumerate}[label=\Alph*.]
\item Rough ER
\item Smooth ER
\item Nucleoli
\item Chloroplasts
\item Centrioles
\end{enumerate}
\textbf{Answer:} B \\
\textbf{Explanation:} Smooth ER contains enzymes involved in the detoxification of drugs and poisons.

\item \textbf{Which of the following is a function of the Golgi apparatus?}
\begin{enumerate}[label=\Alph*.]
\item Calcium storage
\item Protein synthesis
\item Modification of glycoproteins
\item ATP generation
\item DNA replication
\end{enumerate}
\textbf{Answer:} C \\
\textbf{Explanation:} The Golgi modifies products of the ER, such as glycoproteins, as they move from the cis to the trans face.

\item \textbf{The "cis" face of the Golgi apparatus is usually located near the:}
\begin{enumerate}[label=\Alph*.]
\item Plasma membrane
\item Nucleus
\item Endoplasmic reticulum
\item Lysosome
\item Mitochondrion
\end{enumerate}
\textbf{Answer:} C \\
\textbf{Explanation:} The cis face is the "receiving" side that typically faces the ER to receive transport vesicles.

\item \textbf{What is the primary function of lysosomes?}
\begin{enumerate}[label=\Alph*.]
\item Lipid synthesis
\item Intracellular digestion
\item Photosynthesis
\item Storage of cell sap
\item Protein folding
\end{enumerate}
\textbf{Answer:} B \\
\textbf{Explanation:} Lysosomes contain hydrolytic enzymes that break down macromolecules, damaged organelles, and ingested food.

\item \textbf{Tay-Sachs disease is caused by a defect in which organelle?}
\begin{enumerate}[label=\Alph*.]
\item Mitochondrion
\item Peroxisome
\item Lysosome
\item Ribosome
\item Golgi apparatus
\end{enumerate}
\textbf{Answer:} C \\
\textbf{Explanation:} It is a lysosomal storage disease where a lipid-digesting enzyme is missing, leading to brain impairment.

\item \textbf{Which organelle is often the largest in a mature plant cell and holds inorganic ions?}
\begin{enumerate}[label=\Alph*.]
\item Nucleus
\item Chloroplast
\item Central vacuole
\item Mitochondrion
\item Peroxisome
\end{enumerate}
\textbf{Answer:} C \\
\textbf{Explanation:} The central vacuole stores water and ions and plays a major role in plant cell growth.

\item \textbf{Mitochondria are the sites of \underline{\hspace{1cm}}, while chloroplasts are the sites of \underline{\hspace{1cm}}.}
\begin{enumerate}[label=\Alph*.]
\item Photosynthesis; cellular respiration
\item Protein synthesis; lipid synthesis
\item Cellular respiration; photosynthesis
\item Detoxification; digestion
\item Translation; transcription
\end{enumerate}
\textbf{Answer:} C \\
\textbf{Explanation:} Mitochondria generate ATP through respiration; chloroplasts convert solar energy to chemical energy.

\item \textbf{The endosymbiont theory suggests that mitochondria evolved from:}
\begin{enumerate}[label=\Alph*.]
\item Photosynthetic prokaryotes
\item Oxygen-using nonphotosynthetic prokaryotes
\item Folded sections of the plasma membrane
\item Viral infections
\item Ancestral eukaryotic cells
\end{enumerate}
\textbf{Answer:} B \\
\textbf{Explanation:} It is proposed that an early ancestor of eukaryotes engulfed an oxygen-using nonphotosynthetic prokaryote.

\item \textbf{In a mitochondrion, the infoldings of the inner membrane are called:}
\begin{enumerate}[label=\Alph*.]
\item Matrix
\item Stroma
\item Thylakoids
\item Cristae
\item Grana
\end{enumerate}
\textbf{Answer:} D \\
\textbf{Explanation:} Cristae increase the surface area of the inner mitochondrial membrane, enhancing ATP production.

\item \textbf{The fluid found inside a chloroplast but outside the thylakoids is the:}
\begin{enumerate}[label=\Alph*.]
\item Matrix
\item Lumen
\item Cytosol
\item Stroma
\item Cell sap
\end{enumerate}
\textbf{Answer:} D \\
\textbf{Explanation:} The stroma contains the chloroplast DNA, ribosomes, and enzymes needed for sugar synthesis.

\item \textbf{Which organelle produces hydrogen peroxide as a by-product of metabolic reactions?}
\begin{enumerate}[label=\Alph*.]
\item Lysosome
\item Peroxisome
\item Vacuole
\item Golgi
\item Ribosome
\end{enumerate}
\textbf{Answer:} B \\
\textbf{Explanation:} Peroxisomes contain enzymes that transfer hydrogen to oxygen, creating $H_2O_2$, which is then converted to water.

\item \textbf{Which component of the cytoskeleton is the thickest?}
\begin{enumerate}[label=\Alph*.]
\item Microfilaments
\item Intermediate filaments
\item Microtubules
\item Actin filaments
\item Keratin fibers
\end{enumerate}
\textbf{Answer:} C \\
\textbf{Explanation:} Microtubules are hollow rods with a diameter of about 25 nm, the largest of the three fiber types.

\item \textbf{Microtubules are constructed from dimers of:}
\begin{enumerate}[label=\Alph*.]
\item Actin
\item Myosin
\item Keratin
\item Tubulin
\item Collagen
\end{enumerate}
\textbf{Answer:} D \\
\textbf{Explanation:} Tubulin dimers consist of $\alpha$-tubulin and $\beta$-tubulin.

\item \textbf{In animal cells, microtubules grow out from a region near the nucleus called the:}
\begin{enumerate}[label=\Alph*.]
\item Nucleolus
\item Centrosome
\item Basal body
\item Cortex
\item Thylakoid
\end{enumerate}
\textbf{Answer:} B \\
\textbf{Explanation:} The centrosome is the microtubule-organizing center in animal cells, often containing a pair of centrioles.

\item \textbf{What is the arrangement of microtubules in a motile cilium?}
\begin{enumerate}[label=\Alph*.]
\item 9+0
\item 9+2
\item 3+9
\item 5+2
\item 9 triplet arrangement
\end{enumerate}
\textbf{Answer:} B \\
\textbf{Explanation:} Motile cilia and flagella have nine doublets of microtubules in a ring surrounding two central single microtubules.

\item \textbf{Which motor protein is responsible for the bending of cilia and flagella?}
\begin{enumerate}[label=\Alph*.]
\item Myosin
\item Actin
\item Dynein
\item Kinesin
\item Integrin
\item 
\end{enumerate}
\textbf{Answer:} C \\
\textbf{Explanation:} Dynein "feet" walk along adjacent microtubule doublets to cause bending.

\item \textbf{Microfilaments are primarily composed of:}
\begin{enumerate}[label=\Alph*.]
\item Tubulin
\item Keratin
\item Collagen
\item Actin
\item Fibronectin
\end{enumerate}
\textbf{Answer:} D \\
\textbf{Explanation:} Microfilaments are also called actin filaments.

\item \textbf{Which of the following is involved in amoeboid movement?}
\begin{enumerate}[label=\Alph*.]
\item Microtubules and dynein
\item Intermediate filaments and keratin
\item Microfilaments and myosin
\item Central vacuoles and turgor pressure
\item Collagen and integrins
\end{enumerate}
\textbf{Answer:} C \\
\textbf{Explanation:} Localized contractions of actin and myosin drive pseudopodia extension and crawling.

\item \textbf{Intermediate filaments are notable because they are:}
\begin{enumerate}[label=\Alph*.]
\item More permanent than other fibers
\item Composed of tubulin dimers
\item Involved in cytoplasmic streaming
\item Found only in plant cells
\item Used for whole-cell swimming
\end{enumerate}
\textbf{Answer:} A \\
\textbf{Explanation:} Unlike microtubules/microfilaments, intermediate filaments are sturdy and persist even after cell death (e.g., keratin in skin).

\item \textbf{The primary cell wall of plants is composed mostly of:}
\begin{enumerate}[label=\Alph*.]
\item Chitin
\item Peptidoglycan
\item Cellulose
\item Collagen
\item Pectin
\end{enumerate}
\textbf{Answer:} C \\
\textbf{Explanation:} Cellulose microfibrils are the structural component of plant cell walls.

\item \textbf{The thin, glue-like layer between primary walls of adjacent plant cells is the:}
\begin{enumerate}[label=\Alph*.]
\item Secondary wall
\item Plasma membrane
\item Middle lamella
\item Plasmodesma
\item ECM
\end{enumerate}
\textbf{Answer:} C \\
\textbf{Explanation:} The middle lamella is rich in pectins and holds adjacent plant cells together.

\item \textbf{What is the most abundant glycoprotein in the extracellular matrix (ECM) of animal cells?}
\begin{enumerate}[label=\Alph*.]
\item Integrin
\item Fibronectin
\item Collagen
\item Actin
\item Tubulin
\end{enumerate}
\textbf{Answer:} C \\
\textbf{Explanation:} Collagen accounts for about 40\% of the total protein in the human body.

\item \textbf{Cell surface receptor proteins that span the membrane and bind to the ECM are called:}
\begin{enumerate}[label=\Alph*.]
\item Proteoglycans
\item Fibronectins
\item Integrins
\item Desmosomes
\item Aquaporins
\end{enumerate}
\textbf{Answer:} C \\
\textbf{Explanation:} Integrins integrate signals between the ECM and the cytoskeleton.

\item \textbf{Which cell junction is most similar in function to the plasmodesmata of plants?}
\begin{enumerate}[label=\Alph*.]
\item Tight junctions
\item Desmosomes
\item Gap junctions
\item Basal bodies
\item Centrosomes
\end{enumerate}
\textbf{Answer:} C \\
\textbf{Explanation:} Gap junctions provide cytoplasmic channels between adjacent animal cells, just as plasmodesmata do in plants.

\item \textbf{Tight junctions are particularly important for:}
\begin{enumerate}[label=\Alph*.]
\item Fastening muscle cells together
\item Allowing the flow of ions
\item Preventing leakage of extracellular fluid
\item Anchoring the nucleus
\item Photosynthesis
\end{enumerate}
\textbf{Answer:} C \\
\textbf{Explanation:} Tight junctions form continuous seals around cells to prevent leakage across epithelial layers.

\item \textbf{Which junction acts like a "rivet," fastening cells into strong sheets?}
\begin{enumerate}[label=\Alph*.]
\item Tight junction
\item Desmosome
\item Gap junction
\item Plasmodesma
\item Nuclear pore
\end{enumerate}
\textbf{Answer:} B \\
\textbf{Explanation:} Desmosomes (anchoring junctions) fasten cells together, supported by intermediate filaments.

\item \textbf{A protein intended for secretion from the cell is likely synthesized by:}
\begin{enumerate}[label=\Alph*.]
\item Free ribosomes in the cytosol
\item Ribosomes bound to the Rough ER
\item Ribosomes inside mitochondria
\item Ribosomes inside chloroplasts
\item The Golgi apparatus directly
\end{enumerate}
\textbf{Answer:} B \\
\textbf{Explanation:} Bound ribosomes make proteins destined for membranes, packaging in organelles, or export.

\item \textbf{The process of self-digestion by a cell, using its own lysosomes, is called:}
\begin{enumerate}[label=\Alph*.]
\item Phagocytosis
\item Pinocytosis
\item Autophagy
\item Exocytosis
\item Cell fractionation
\end{enumerate}
\textbf{Answer:} C \\
\textbf{Explanation:} Autophagy is the recycling of the cell's own organic material.

\item \textbf{Which of the following would NOT be part of the endomembrane system?}
\begin{enumerate}[label=\Alph*.]
\item Nuclear envelope
\item Golgi apparatus
\item Mitochondrion
\item Plasma membrane
\item Lysosome
\end{enumerate}
\textbf{Answer:} C \\
\textbf{Explanation:} Mitochondria and chloroplasts are not part of the endomembrane system; they have double membranes and different origins.

\item \textbf{What is the function of the smooth ER in muscle cells?}
\begin{enumerate}[label=\Alph*.]
\item Protein synthesis
\item Sorting proteins for secretion
\item Storage and release of calcium ions
\item DNA replication
\item Photosynthesis
\end{enumerate}
\textbf{Answer:} C \\
\textbf{Explanation:} In muscle cells, a specialized smooth ER (sarcoplasmic reticulum) regulates calcium levels for contraction.

\item \textbf{Grana, thylakoids, and stroma are all components found in:}
\begin{enumerate}[label=\Alph*.]
\item Vacuoles
\item Chloroplasts
\item Mitochondria
\item Nuclei
\item Lysosomes
\end{enumerate}
\textbf{Answer:} B \\
\textbf{Explanation:} These are the internal structural components of chloroplasts.

\item \textbf{The core of a microvillus is composed of:}
\begin{enumerate}[label=\Alph*.]
\item Microtubules
\item Intermediate filaments
\item Microfilaments
\item Myosin
\item Keratin
\end{enumerate}
\textbf{Answer:} C \\
\textbf{Explanation:} Bundles of microfilaments (actin) provide the structural core for microvilli.

\item \textbf{Which tool would you use to see the changes in shape of a \textit{living} white blood cell?}
\begin{enumerate}[label=\Alph*.]
\item SEM
\item TEM
\item Light Microscope
\item Cryo-EM
\item X-ray crystallography
\end{enumerate}
\textbf{Answer:} C \\
\textbf{Explanation:} Only light microscopy allows for the observation of living specimens.

\item \textbf{The small circular spots seen on a nuclear lamina TEM are:}
\begin{enumerate}[label=\Alph*.]
\item Ribosomes
\item Nucleoli
\item Nuclear pores
\item Histones
\item Chromatin loops
\end{enumerate}
\textbf{Answer:} C \\
\textbf{Explanation:} Nuclear pores perforate the envelope and allow transport between the nucleus and cytoplasm.

\item \textbf{Which organelle is responsible for synthesizing oils, steroids, and phospholipids?}
\begin{enumerate}[label=\Alph*.]
\item Rough ER
\item Smooth ER
\item Golgi
\item Lysosome
\item Peroxisome
\end{enumerate}
\textbf{Answer:} B \\
\textbf{Explanation:} Smooth ER is the primary site of lipid synthesis.

\item \textbf{What happens to a food vacuole after it is formed by phagocytosis?}
\begin{enumerate}[label=\Alph*.]
\item It fuses with the Golgi
\item It fuses with a lysosome
\item It is expelled immediately
\item It becomes a mitochondrion
\item It enters the nucleus
\end{enumerate}
\textbf{Answer:} B \\
\textbf{Explanation:} Lysosomes fuse with food vacuoles to digest the contents using hydrolytic enzymes.

\item \textbf{Cilia move in a/an \underline{\hspace{1cm}} motion, while flagella move in a/an \underline{\hspace{1cm}} motion.}
\begin{enumerate}[label=\Alph*.]
\item Undulating; back-and-forth
\item Back-and-forth; undulating
\item Circular; linear
\item Rotating; vibrating
\item Stiff; limp
\end{enumerate}
\textbf{Answer:} B \\
\textbf{Explanation:} Cilia have power and recovery strokes; flagella undulate like a fish tail.

\item \textbf{A "primary cilium" differs from motile cilia because it:}
\begin{enumerate}[label=\Alph*.]
\item Is longer
\item Acts as a sensory antenna
\item Has a 9+2 pattern
\item Is used for swimming
\item Lacks a basal body
\end{enumerate}
\textbf{Answer:} B \\
\textbf{Explanation:} Primary cilia are nonmotile and transmit molecular signals from the environment.

\item \textbf{Which protein subunits make up the fibrous subunits of intermediate filaments?}
\begin{enumerate}[label=\Alph*.]
\item Tubulins
\item Actins
\item Keratins
\item Myosins
\item Dyneins
\end{enumerate}
\textbf{Answer:} C \\
\textbf{Explanation:} Intermediate filaments are constructed from various proteins, most commonly keratins.

\item \textbf{Cytoplasmic streaming in plants is driven by:}
\begin{enumerate}[label=\Alph*.]
\item Microtubules and dynein
\item Actin-myosin interactions
\item Turgor pressure only
\item The Golgi apparatus
\item Flagellar rotation
\end{enumerate}
\textbf{Answer:} B \\
\textbf{Explanation:} Myosin motors attached to organelles move along actin tracks, driving the cytosol flow.

\item \textbf{Which structure connects the cytoplasms of adjacent plant cells?}
\begin{enumerate}[label=\Alph*.]
\item Gap junction
\item Tight junction
\item Desmosome
\item Plasmodesma
\item Middle lamella
\end{enumerate}
\textbf{Answer:} D \\
\textbf{Explanation:} Plasmodesmata are membrane-lined channels through plant cell walls.

\item \textbf{In the endosymbiont theory, the ancestral photosynthetic prokaryote became the:}
\begin{enumerate}[label=\Alph*.]
\item Mitochondrion
\item Chloroplast
\item Nucleus
\item Ribosome
\item Peroxisome
\end{enumerate}
\textbf{Answer:} B \\
\textbf{Explanation:} Engulfed photosynthetic prokaryotes are the proposed ancestors of chloroplasts.

\item \textbf{Mitochondria and chloroplasts are "autonomous" because they:}
\begin{enumerate}[label=\Alph*.]
\item Do not need the cell to survive
\item Have their own DNA and ribosomes
\item Can leave the cell at will
\item Are made of protein only
\item Lacking membranes
\end{enumerate}
\textbf{Answer:} B \\
\textbf{Explanation:} They contain circular DNA and ribosomes and can reproduce somewhat independently within the cell.

\item \textbf{What is the function of glyoxysomes in plant seeds?}
\begin{enumerate}[label=\Alph*.]
\item Photosynthesis
\item Starch storage
\item Converting fatty acids to sugar
\item Protein synthesis
\item Detoxifying alcohol
\end{enumerate}
\textbf{Answer:} C \\
\textbf{Explanation:} Glyoxysomes are specialized peroxisomes that initiate fatty acid conversion for the seedling's energy.

\item \textbf{Proteins that catalyze the first steps of sugar breakdown are usually made by:}
\begin{enumerate}[label=\Alph*.]
\item Bound ribosomes
\item Free ribosomes
\item The Golgi
\item Mitochondria
\item Lysosomes
\end{enumerate}
\textbf{Answer:} B \\
\textbf{Explanation:} Proteins functioning within the cytosol are typically made on free ribosomes.

\item \textbf{A nuclear matrix is a:}
\begin{enumerate}[label=\Alph*.]
\item Liquid inside the nucleus
\item Double membrane
\item Framework of protein fibers in the nucleus
\item Hole in the nuclear envelope
\item Region for ribosome synthesis
\end{enumerate}
\textbf{Answer:} C \\
\textbf{Explanation:} The matrix provides a structural framework throughout the nuclear interior.

\item \textbf{How many chromosomes are in a typical human sex cell?}
\begin{enumerate}[label=\Alph*.]
\item 46
\item 8
\item 23
\item 4
\item 2
\end{enumerate}
\textbf{Answer:} C \\
\textbf{Explanation:} Human somatic cells have 46; sex cells (gametes) have 23.

\item \textbf{Which structure is responsible for the "9+0" arrangement?}
\begin{enumerate}[label=\Alph*.]
\item Motile flagellum
\item Basal body
\item Primary cilium
\item Cytoplasm
\item Both B and C
\end{enumerate}
\textbf{Answer:} E \\
\textbf{Explanation:} Basal bodies and nonmotile primary cilia lack the central pair of microtubules.

\item \textbf{Dynein "walking" requires energy in the form of:}
\begin{enumerate}[label=\Alph*.]
\item Sunlight
\item Glucose
\item ATP
\item NADH
\item Heat
\end{enumerate}
\textbf{Answer:} C \\
\textbf{Explanation:} Motor proteins like dynein are ATP-powered.

\item \textbf{In the "cisternal maturation model," Golgi cisternae:}
\begin{enumerate}[label=\Alph*.]
\item Remain static while vesicles move cargo
\item Disappear and reappear
\item Move from the cis to the trans face
\item Are only found in prokaryotes
\item Merge with the nuclear envelope
\end{enumerate}
\textbf{Answer:} C \\
\textbf{Explanation:} This model proposes that cisternae themselves progress forward, carrying and modifying cargo.

\item \textbf{Which organelle is responsible for the hydrolysis of macromolecules in animal cells?}
\begin{enumerate}[label=\Alph*.]
\item Ribosome
\item Lysosome
\item Smooth ER
\item Centrosome
\item Mitochondrion
\end{enumerate}
\textbf{Answer:} B \\
\textbf{Explanation:} Lysosomes are the primary digestive compartments.

\item \textbf{A cell has reached its maximum size when:}
\begin{enumerate}[label=\Alph*.]
\item It runs out of DNA
\item Its volume exceeds its surface area's capacity to service it
\item It hits the cell wall
\item It has 100 mitochondria
\item It becomes multicellular
\end{enumerate}
\textbf{Answer:} B \\
\textbf{Explanation:} The surface area to volume ratio limits the exchange of nutrients and wastes.

\item \textbf{Which of the following is NOT a parameter of microscopy?}
\begin{enumerate}[label=\Alph*.]
\item Resolution
\item Contrast
\item Magnification
\item Fractionation
\item Clarity
\end{enumerate}
\textbf{Answer:} D \\
\textbf{Explanation:} Fractionation is a biochemical technique, not a microscopy parameter.

\item \textbf{Microscopy + Biochemistry = ?}
\begin{enumerate}[label=\Alph*.]
\item Histology
\item Cytology
\item Cell Biology
\item Genetics
\item Metabolism
\end{enumerate}
\textbf{Answer:} C \\
\textbf{Explanation:} Modern cell biology integrates the study of structure (cytology) and function (biochemistry).

\item \textbf{Contractile vacuoles in freshwater protists function to:}
\begin{enumerate}[label=\Alph*.]
\item Digest food
\item Pump out excess water
\item Store pigments
\item Synthesize proteins
\item Move the cell
\end{enumerate}
\textbf{Answer:} B \\
\textbf{Explanation:} They maintain ion concentrations by expelling water that enters via osmosis.

\item \textbf{What is the function of the "cortex" in a cell?}
\begin{enumerate}[label=\Alph*.]
\item It is the center of the nucleus
\item It is the outer cytoplasmic layer with a gel-like consistency
\item It is the site of ATP production
\item It is the boundary of the cell wall
\item It holds the chromosomes
\end{enumerate}
\textbf{Answer:} B \\
\textbf{Explanation:} The cortex is reinforced by a 3D network of microfilaments.

\item \textbf{The extracellular matrix (ECM) communicates with the nucleus via:}
\begin{enumerate}[label=\Alph*.]
\item Direct DNA contact
\item Integrins and cytoskeletal signaling
\item Pores in the plasma membrane
\item Diffusion of collagen
\item Mitochondria
\end{enumerate}
\textbf{Answer:} B \\
\textbf{Explanation:} Mechanical and chemical signaling pathways link the ECM to the interior.

\item \textbf{Plant cells lack \underline{\hspace{1cm}}, which are common in animal cells.}
\begin{enumerate}[label=\Alph*.]
\item Mitochondria
\item Centrioles
\item Ribosomes
\item Golgi
\item Nuclei
\end{enumerate}
\textbf{Answer:} B \\
\textbf{Explanation:} While plants have centrosomes, they typically lack the centrioles found in animals.

\item \textbf{A proteoglycan molecule is mostly composed of:}
\begin{enumerate}[label=\Alph*.]
\item Amino acids
\item Lipids
\item Carbohydrates
\item DNA
\item Metal ions
\end{enumerate}
\textbf{Answer:} C \\
\textbf{Explanation:} Proteoglycans can be up to 95\% carbohydrate.

\item \textbf{Which of these statements about the cell wall is FALSE?}
\begin{enumerate}[label=\Alph*.]
\item It is much thicker than the plasma membrane.
\item Its composition is identical in all plant species.
\item It prevents excessive water uptake.
\item It maintains cell shape.
\item It can have multiple layers.
\end{enumerate}
\textbf{Answer:} B \\
\textbf{Explanation:} Chemical composition varies between species and even cell types.

\item \textbf{The protein that "zips" together animal cells to prevent fluid movement is:}
\begin{enumerate}[label=\Alph*.]
\item Keratin
\item Collagen
\item Tight junction protein
\item Tubulin
\item Myosin
\end{enumerate}
\textbf{Answer:} C \\
\textbf{Explanation:} Tight junctions create a watertight seal.

\item \textbf{If a cell's surface area is 150 units$^2$ and volume is 125 units$^3$, its SA:V ratio is:}
\begin{enumerate}[label=\Alph*.]
\item 1.2
\item 6
\item 0.83
\item 18750
\item 25
\end{enumerate}
\textbf{Answer:} A \\
\textbf{Explanation:} $150 / 125 = 1.2$.

\item \textbf{An amyloplast is a type of \underline{\hspace{1cm}} that stores \underline{\hspace{1cm}}.}
\begin{enumerate}[label=\Alph*.]
\item Mitochondrion; energy
\item Plastid; starch
\item Lysosome; enzymes
\item Vacuole; water
\item Ribosome; protein
\end{enumerate}
\textbf{Answer:} B \\
\textbf{Explanation:} Amyloplasts are colorless plastids that store starch.

\item \textbf{Mitochondrial DNA is usually:}
\begin{enumerate}[label=\Alph*.]
\item Linear
\item Double-helical but not circular
\item Circular
\item Single-stranded
\item Found only in the intermembrane space
\end{enumerate}
\textbf{Answer:} C \\
\textbf{Explanation:} Like bacterial DNA, mitochondrial DNA is circular.

\item \textbf{In centrifugation, the liquid remaining above the pellet is the:}
\begin{enumerate}[label=\Alph*.]
\item Homogenate
\item Filtrate
\item Supernatant
\item Matrix
\item Solute
\end{enumerate}
\textbf{Answer:} C \\
\textbf{Explanation:} The supernatant is poured into a new tube for higher-speed centrifugation.

\item \textbf{Which microscope would be best for seeing the internal \textit{details} of a ribosome?}
\begin{enumerate}[label=\Alph*.]
\item Standard LM
\item SEM
\item TEM or Cryo-EM
\item Phase-contrast
\item Magnifying glass
\end{enumerate}
\textbf{Answer:} C \\
\textbf{Explanation:} Ribosomes are extremely small (approx 20-30nm), requiring the high resolution of TEM or Cryo-EM.

\item \textbf{"Cell sap" is found in:}
\begin{enumerate}[label=\Alph*.]
\item Cytosol
\item Nucleus
\item Central vacuole
\item Blood
\item Thylakoids
\end{enumerate}
\textbf{Answer:} C \\
\textbf{Explanation:} Cell sap is the solution inside the plant's central vacuole.

\item \textbf{Detoxification in the liver involves adding \underline{\hspace{1cm}} groups to drugs.}
\begin{enumerate}[label=\Alph*.]
\item Methyl
\item Hydroxyl
\item Carboxyl
\item Phosphate
\item Amino
\end{enumerate}
\textbf{Answer:} B \\
\textbf{Explanation:} Adding hydroxyl groups makes the drugs more water-soluble for excretion.

\item \textbf{What is the primary function of the nuclear lamina?}
\begin{enumerate}[label=\Alph*.]
\item DNA replication
\item Ribosome assembly
\item Maintaining nuclear shape
\item RNA transport
\item ATP synthesis
\end{enumerate}
\textbf{Answer:} C \\
\textbf{Explanation:} It provides mechanical support to the nuclear envelope.

\item \textbf{Which structure is responsible for the "power stroke" in cilia?}
\begin{enumerate}[label=\Alph*.]
\item Myosin
\item Dynein
\item Actin
\item Keratin
\item Collagen
\end{enumerate}
\textbf{Answer:} B \\
\textbf{Explanation:} Dynein arms provide the force for the ciliary beat.

\item \textbf{A "microsome" in fractionation is a piece of:}
\begin{enumerate}[label=\Alph*.]
\item DNA
\item Mitochondrion
\item Membrane (plasma or internal)
\item Ribosome
\item Nucleolus
\end{enumerate}
\textbf{Answer:} C \\
\textbf{Explanation:} Microsomes are vesicles formed from disrupted ER and plasma membranes.

\item \textbf{Which of these is NOT a role of the cytoskeleton?}
\begin{enumerate}[label=\Alph*.]
\item Support and shape
\item Cell motility
\item Organelle anchorage
\item DNA synthesis
\item Signal transmission
\end{enumerate}
\textbf{Answer:} D \\
\textbf{Explanation:} DNA synthesis occurs in the nucleus via enzymes, not directly by the cytoskeleton.

\item \textbf{Integrins bind to \underline{\hspace{1cm}} on the outside and \underline{\hspace{1cm}} on the inside.}
\begin{enumerate}[label=\Alph*.]
\item DNA; RNA
\item Collagen; Myosin
\item Fibronectin; Microfilaments
\item Water; Ions
\item Sugar; Starch
\end{enumerate}
\textbf{Answer:} C \\
\textbf{Explanation:} This linkage allows the ECM to influence cell behavior.

\item \textbf{What occurs in the "transitional ER"?}
\begin{enumerate}[label=\Alph*.]
\item Lipids are synthesized
\item Transport vesicles bud off
\item Ribosomes are assembled
\item DNA is stored
\item Proteins are degraded
\end{enumerate}
\textbf{Answer:} B \\
\textbf{Explanation:} Vesicles bud from this specialized region of the rough ER.

\item \textbf{Most secretory proteins are \underline{\hspace{1cm}}, which have carbohydrates attached.}
\begin{enumerate}[label=\Alph*.]
\item Lipoproteins
\item Nucleoproteins
\item Glycoproteins
\item Phosphoproteins
\item Metalloproteins
\end{enumerate}
\textbf{Answer:} C \\
\textbf{Explanation:} Carbohydrates are added to proteins in the ER lumen.

\item \textbf{The theory that a host cell and engulfed cell merged into a single organism is:}
\begin{enumerate}[label=\Alph*.]
\item Cell theory
\item Germ theory
\item Endosymbiont theory
\item Theory of relativity
\item Evolutionary synthesis
\end{enumerate}
\textbf{Answer:} C \\
\textbf{Explanation:} This explains the origin of mitochondria and chloroplasts.

\item \textbf{Chromatin consists of DNA and:}
\begin{enumerate}[label=\Alph*.]
\item RNA only
\item Lipids
\item Proteins (like histones)
\item Carbohydrates
\item ATP
\end{enumerate}
\textbf{Answer:} C \\
\textbf{Explanation:} Proteins help coil the DNA to fit into the nucleus.

\item \textbf{Which organelle is a "warehouse" for receiving, sorting, and shipping?}
\begin{enumerate}[label=\Alph*.]
\item ER
\item Golgi apparatus
\item Nucleus
\item Lysosome
\item Vacuole
\end{enumerate}
\textbf{Answer:} B \\
\textbf{Explanation:} The Golgi is the shipping and receiving center.

\item \textbf{What is the function of pectins in the middle lamella?}
\begin{enumerate}[label=\Alph*.]
\item Energy storage
\item DNA protection
\item Gluing cells together
\item Speeding up reactions
\item Transporting water
\end{enumerate}
\textbf{Answer:} C \\
\textbf{Explanation:} Pectins are sticky polysaccharides that act as glue.

\item \textbf{A "mesosome" is sometimes found in prokaryotes; what does it resemble?}
\begin{enumerate}[label=\Alph*.]
\item A nucleus
\item An infolding of the plasma membrane
\item A ribosome
\item A cell wall
\item A flagellum
\end{enumerate}
\textbf{Answer:} B \\
\textbf{Explanation:} Although once thought to be organelles, they are often artifacts of preparation.

\item \textbf{Intermediate filaments are made of \underline{\hspace{1cm}} proteins.}
\begin{enumerate}[label=\Alph*.]
\item Globular
\item Fibrous
\item Small
\item Circular
\item Temporary
\end{enumerate}
\textbf{Answer:} B \\
\textbf{Explanation:} They are fibrous proteins coiled into cables.

\item \textbf{Which of the following is a "communicating junction"?}
\begin{enumerate}[label=\Alph*.]
\item Tight junction
\item Desmosome
\item Gap junction
\item Basal body
\item Cell wall
\end{enumerate}
\textbf{Answer:} C \\
\textbf{Explanation:} Gap junctions allow for communication between cells.

\item \textbf{Eukaryotic cells are typically \underline{\hspace{1cm}} times larger than bacteria.}
\begin{enumerate}[label=\Alph*.]
\item 2
\item 10-100
\item 1000
\item 0.5
\item 1,000,000
\end{enumerate}
\textbf{Answer:} B \\
\textbf{Explanation:} Typical bacteria are 1-5 $\mu$m; eukaryotes are 10-100 $\mu$m.

\item \textbf{Which is NOT a function of the central vacuole?}
\begin{enumerate}[label=\Alph*.]
\item Storage of ions
\item Waste disposal
\item ATP synthesis
\item Cell growth
\item Protection against herbivores
\end{enumerate}
\textbf{Answer:} C \\
\textbf{Explanation:} Mitochondria and chloroplasts synthesize ATP, not the vacuole.

\item \textbf{A "granum" is a stack of:}
\begin{enumerate}[label=\Alph*.]
\item Cristae
\item Thylakoids
\item Ribosomes
\item Nuclei
\item Golgi sacs
\end{enumerate}
\textbf{Answer:} B \\
\textbf{Explanation:} Grana are stacks of thylakoids inside chloroplasts.

\item \textbf{The name "Endoplasmic Reticulum" means:}
\begin{enumerate}[label=\Alph*.]
\item Little net within the cytoplasm
\item Powerhouse of the cell
\item Brain of the cell
\item Shipping container
\item Acidic sac
\end{enumerate}
\textbf{Answer:} A \\
\textbf{Explanation:} \textit{Endo} (within), \textit{plasmic} (cytoplasm), \textit{reticulum} (little net).

\item \textbf{Which microscopy breakthrough allows for breaking the 200nm resolution barrier?}
\begin{enumerate}[label=\Alph*.]
\item Phase-contrast
\item Super-resolution microscopy
\item Brightfield
\item SEM
\item Standard LM
\end{enumerate}
\textbf{Answer:} B \\
\textbf{Explanation:} Super-resolution techniques can distinguish structures as small as 10-20 nm.

\item \textbf{The cell is the \underline{\hspace{1cm}} unit of structure and function.}
\begin{enumerate}[label=\Alph*.]
\item Largest
\item Only
\item Basic
\item Most complex
\item Secondary
\end{enumerate}
\textbf{Answer:} C \\
\textbf{Explanation:} All living things are made of cells, the basic unit of life.

\item \textbf{Which organelle's number correlates with a cell's metabolic activity?}
\begin{enumerate}[label=\Alph*.]
\item Nucleus
\item Mitochondrion
\item Vacuole
\item Cell wall
\item Nucleolus
\end{enumerate}
\textbf{Answer:} B \\
\textbf{Explanation:} Active cells (like muscle cells) have many more mitochondria.

\end{enumerate}
\end{document}